\documentclass[10pt,a4paper]{amsart}
\usepackage[margin=19mm]{geometry}
\usepackage{amsmath,amssymb,mathtools}
\usepackage[scr=rsfs,cal=euler]{mathalfa}
\usepackage{xfrac}
\usepackage[unicode,hidelinks,bookmarks=false]{hyperref}
\newcommand{\ie}{i.e.\ }
\newcommand{\cf}{cf.\ }
\newcommand{\resp}{respectively\ }
\newcommand{\Subsection}{\subsection}
\providecommand{\nobreakdash}{\nobreak\hbox{-}}
\setlength{\emergencystretch}{2em}
\multlinegap0pt
\usepackage{bm}
\usepackage{bbm}
\usepackage{dsfont}
\usepackage{xspace}
\usepackage{cancel}
\usepackage{mathtools}
\usepackage{amsthm}
\makeatletter
\@ifundefined{theorem}{\newtheorem{theorem}{Theorem}}{}
\makeatother
\newcommand{\R}{\mathbb{R}}
\newcommand{\Z}{\mathbb{Z}}
\title[The Exponent Set: A Third Natural Extension of the Mandelbro]{The Exponent Set: A Third Natural Extension of the Mandelbrot-Julia Framework}
\author{Yuchen Brian Shen}
\date{Source: arXiv:2608.07560v1 (CC BY 4.0)}
\begin{document}
\maketitle

\noindent\emph{Source and licence.} The Exponent Set: A Third Natural Extension of the Mandelbrot-Julia Framework. Version arXiv:2608.07560v1 (CC BY 4.0), distributed under the Creative Commons Attribution 4.0 International licence, as recorded in the arXiv licence metadata for this exact version and shown on the abstract page https://arxiv.org/abs/2608.07560v1. Full rights notice: licenses/arxiv-2608.07560-CC-BY-4.0.txt.\par\smallskip

\noindent\textit{Editorial scope.} Counted unit: the Theorem whose source label is \texttt{thm:no-universal-radius-negative-real} in the article (source0.tex lines 741--757, source label \texttt{thm:no-universal-radius-negative-real}), delivered together with the article's own complete original proof of it (source0.tex lines 759--1147, one \texttt{proof} environment of 389 lines). No mathematics is written, repaired or re-derived: statement and proof are the article's own text line for line. Required context retained uncounted: none. Every definition, display and label of those lines is retained unchanged. Omitted from this package and not redistributed: 1--740, 758--758, 1148--2522. Nothing inside a retained excerpt is omitted or altered: every retained body is delivered exactly as the article's lines are written, so no omission receipt of any kind accompanies this package. Citations: the retained excerpts carry no citation command at all, so no bibliography entry of the article is transcribed and no reference is redistributed. Reference adaptation: none was needed and none was made; every reference inside the delivered unit resolves inside it, so no unresolved reference or citation is produced. Printed slips kept literal: no printed word, symbol, hypothesis or number of the counted statement or of its proof was altered. Layout and class: the article's own class is \texttt{amsart} with packages including amsmath, amssymb, amsthm, array, graphicx, placeins, flafter, natbib, xcolor, float, hyperref; the wrapper is the lane's pinned \texttt{amsart} editorial wrapper at 10pt on A4, which restates the theorem-like environments the retained text needs and adds the article's own macro definitions that the retained text needs (\texttt{\textbackslash R}, \texttt{\textbackslash Z}), with extra wrapper packages (bm, bbm, dsfont, xspace, cancel, mathtools). The delivered numbering is the unit's own. Assets: no retained span contains a figure environment, a graphics-inclusion command, a PDF-page inclusion, a table or a TikZ picture; no image file is copied into this package, the package declares no asset dependency and no image file is redistributed with it. Licence: version arXiv:2608.07560v1 of this article is distributed under the Creative Commons Attribution 4.0 International licence, as recorded in the arXiv licence metadata for this exact version and shown on the abstract page https://arxiv.org/abs/2608.07560v1. A full rights notice is delivered at licenses/arxiv-2608.07560-CC-BY-4.0.txt. Characters: every retained excerpt is pure ASCII, so the delivered engine, XeLaTeX with its default Latin Modern text font, reports no missing character.
\begin{theorem}[Absence of a Universal Escape Radius for Negative Real Initial Values]
\label{thm:no-universal-radius-negative-real}
Let $c=0$ and let $z_0\in\R$ satisfy
\[
z_0<0,
\qquad
z_0\ne -1.
\]
Then the Exponent Set $E(z_0,0)$ admits no finite universal escape radius.

More precisely, for every $r>0$, there exists an exponent $x\in E(z_0,0)$
such that the orbit generated by
\[
z_{n+1}=z_n^{\,x}
\]
satisfies $|z_n|>r$ for some $n\ge 0$, while the orbit remains bounded.
\end{theorem}

\begin{proof}
Write
\[
z_0=-\rho,
\qquad
\rho>0.
\]
Since $z_0\ne -1$, we have $\rho\ne 1$. Define
\[
a=\ln(\rho).
\]
Then
\[
a\in\R\setminus\{0\}.
\]
By the principal-argument convention, the principal argument of any negative
real number is $\pi$. Therefore
\[
\log(z_0)=a+i\pi.
\]
Here $\pi$, not $-\pi$, is the principal argument of $z_0<0$.

Let
\[
A=|a|.
\]
Then $A>0$. Define
\[
\sigma=-\frac{a}{A}.
\]
Then $\sigma\in\{-1,1\}$ and
\[
\sigma A=-a.
\]

Fix $r>0$. We will construct an exponent $x$ such that the orbit satisfies
\[
z_0\mapsto z_1\mapsto 1\mapsto 1\mapsto \cdots
\]
and such that $|z_1|>r$.

For each integer $N\ge 1$ and each
\[
u\in[\pi/2,3\pi/4],
\]
define
\[
S_N(u)=u+\pi N,
\]
\[
B_N(u)=\frac{2\pi S_N(u)}{A},
\]
and
\[
C_N(u)=u(u+2\pi N).
\]
Now define
\[
M_N(u)
=
\frac{B_N(u)+\sqrt{B_N(u)^2+4C_N(u)}}{2}.
\]
Then $M_N(u)>0$, and $M_N(u)$ is the positive root of
\[
M^2-B_N(u)M-C_N(u)=0.
\]
Equivalently,
\[
M_N(u)^2-C_N(u)=B_N(u)M_N(u).
\]

Define
\[
G_N(u)=\frac{M_N(u)S_N(u)}{\pi A}.
\]
The displayed radical formula shows that $M_N$ is continuously
differentiable on $[\pi/2,3\pi/4]$. Consequently, $G_N$ is continuously
differentiable there as well.

We first show that, for some large $N$, the function $G_N$ takes an integer
value. Since $B_N(u)>0$ and $C_N(u)>0$, we have
\[
M_N(u)>B_N(u)
\]
for all $u\in[\pi/2,3\pi/4]$.

Moreover, differentiating the equation
\[
M_N(u)^2-B_N(u)M_N(u)-C_N(u)=0
\]
gives
\[
(2M_N(u)-B_N(u))M_N'(u)=B_N'(u)M_N(u)+C_N'(u).
\]
Here
\[
2M_N(u)-B_N(u)=\sqrt{B_N(u)^2+4C_N(u)}>0,
\]
while $B_N'(u)>0$ and $C_N'(u)>0$. Hence
\[
M_N'(u)>0.
\]
Therefore
\[
G_N'(u)
=
\frac{M_N'(u)S_N(u)+M_N(u)}{\pi A}
>
\frac{M_N(u)}{\pi A}
>
\frac{B_N(u)}{\pi A}
=
\frac{2S_N(u)}{A^2}.
\]
It follows that
\[
G_N(3\pi/4)-G_N(\pi/2)
>
\int_{\pi/2}^{3\pi/4}\frac{2(u+\pi N)}{A^2}\,du.
\]
The right-hand side tends to $+\infty$ as $N\to\infty$. Hence we may choose
$N$ sufficiently large such that
\[
G_N(3\pi/4)-G_N(\pi/2)>1.
\]
We also choose $N$ sufficiently large so that
\[
M_N(u)>\max\{0,\ln(r)\}
\]
for every $u\in[\pi/2,3\pi/4]$. This is possible because
\[
M_N(u)>B_N(u)=\frac{2\pi(u+\pi N)}{A}\to+\infty
\]
uniformly for $u\in[\pi/2,3\pi/4]$.

Define
\[
m=\left\lceil G_N(\pi/2)\right\rceil.
\]
Because $G_N(3\pi/4)-G_N(\pi/2)>1$ and both endpoint values are positive, we
have
\[
1\le m\le G_N(3\pi/4).
\]
By the intermediate value theorem, there exists
\[
u\in[\pi/2,3\pi/4]
\]
such that
\[
G_N(u)=m.
\]

Now set
\[
M=M_N(u),
\qquad
S=S_N(u),
\]
and define
\[
\theta=\sigma u,
\qquad
p=\sigma N,
\qquad
T=\theta+2\pi p.
\]
Since $\sigma=\pm 1$, we have $p\in\Z$. Also,
\[
T=\sigma(u+2\pi N).
\]
Because $u\in[\pi/2,3\pi/4]$, we have
\[
\theta\in[-3\pi/4,-\pi/2]\cup[\pi/2,3\pi/4]\subset(-\pi,\pi].
\]
Thus $\theta$ lies in the principal argument range.

Finally, define
\[
x=\frac{M+iT}{a+i\pi}.
\]
This is well-defined because $a+i\pi\ne 0$.

We now compute the orbit. Since $\log(z_0)=a+i\pi$, the first iterate is
\[
z_1
=
z_0^{\,x}
=
\exp(x\log(z_0))
=
\exp(M+iT).
\]
Since $T=\theta+2\pi p$ and $p\in\Z$, this becomes
\[
z_1=e^M e^{i\theta}.
\]
Therefore
\[
|z_1|=e^M.
\]
By the choice of $N$, we have $M>\max\{0,\ln(r)\}$. Hence
\[
|z_1|=e^M>r.
\]

It remains to prove that the orbit is bounded. Since $\theta\in(-\pi,\pi]$,
the principal-value logarithm of $z_1$ is
\[
\log(z_1)=M+i\theta.
\]
Thus
\[
z_2
=
z_1^{\,x}
=
\exp(x\log(z_1))
=
\exp\left(\frac{M+iT}{a+i\pi}(M+i\theta)\right).
\]

We now compute the exponent explicitly. Since
\[
T=\sigma(u+2\pi N)
\qquad\text{and}\qquad
\theta=\sigma u,
\]
we have
\[
T\theta
=
\sigma(u+2\pi N)\sigma u
=
u(u+2\pi N)
=
C_N(u).
\]
Also,
\[
T+\theta
=
\sigma(u+2\pi N)+\sigma u
=
2\sigma(u+\pi N)
=
2\sigma S.
\]

Since $M=M_N(u)$ is the positive root of
\[
M^2-B_N(u)M-C_N(u)=0,
\]
we have
\[
M^2-C_N(u)=B_N(u)M.
\]
Using $T\theta=C_N(u)$, this gives
\[
M^2-T\theta=B_N(u)M.
\]
Because
\[
B_N(u)=\frac{2\pi S}{A},
\]
we obtain
\[
M^2-T\theta
=
\frac{2\pi SM}{A}.
\]
On the other hand, from
\[
G_N(u)=\frac{M_N(u)S_N(u)}{\pi A}=m
\]
and the definitions $M=M_N(u)$ and $S=S_N(u)$, we get
\[
MS=\pi A m.
\]
Therefore
\[
M^2-T\theta
=
\frac{2\pi MS}{A}
=
2\pi^2m.
\]

Similarly, since $T+\theta=2\sigma S$, we have
\[
M(T+\theta)
=
2\sigma MS.
\]
Using $MS=\pi A m$, this becomes
\[
M(T+\theta)
=
2\sigma\pi A m.
\]
Since $\sigma A=-a$, we obtain
\[
M(T+\theta)
=
-2\pi a m.
\]

Therefore
\[
\begin{aligned}
(M+iT)(M+i\theta)
&=
M^2-T\theta+iM(T+\theta) \\
&=
2\pi^2m-2\pi i a m.
\end{aligned}
\]
The last expression factors as
\[
2\pi^2m-2\pi i a m
=
-2\pi i m(a+i\pi).
\]
Hence
\[
(M+iT)(M+i\theta)
=
-2\pi i m(a+i\pi).
\]
Since
\[
x=\frac{M+iT}{a+i\pi}
\qquad\text{and}\qquad
\log(z_1)=M+i\theta,
\]
we conclude that
\[
x\log(z_1)
=
\frac{M+iT}{a+i\pi}(M+i\theta)
=
-2\pi i m.
\]
Therefore
\[
z_2
=
\exp(x\log(z_1))
=
\exp(-2\pi i m)
=
1.
\]

Now, using the principal-value logarithm, $\log(1)=0$. Hence
\[
z_3
=
1^x
=
\exp(x\log(1))
=
\exp(0)
=
1.
\]
By induction,
\[
z_n=1
\qquad\text{for all }n\ge 2.
\]
Thus the orbit is
\[
z_0,\quad z_1,\quad 1,\quad 1,\quad 1,\ldots
\]
and is bounded by
\[
B=\max\{|z_0|,|z_1|,1\}<\infty.
\]
All iterates in this constructed orbit are nonzero, so the orbit is
well-defined for every $n\ge 0$. Therefore
\[
x\in E(z_0,0).
\]

Since for the arbitrary radius $r>0$ we have constructed an exponent
$x\in E(z_0,0)$ whose orbit satisfies $|z_1|>r$, no finite number $r$ can
serve as a universal escape radius for $E(z_0,0)$.
\end{proof}

\end{document}
